\documentclass[10pt]{article}
\usepackage{amsmath,amsthm,amssymb,amsfonts}
\usepackage[margin=1in]{geometry}
\usepackage{hyperref}
\usepackage{url}
\newtheorem{theorem}{Theorem}
\title{Supplementary Note: Full Algebraic Verification of the Almost Orthogonal Property}
\author{Kazufumi Ito \and Tong Wu \and Jun Zou}
\date{}
\begin{document}
\maketitle

\begin{abstract}
This note provides the fully expanded algebraic verification of the almost
orthogonal property used in the Appendix of the manuscript
``\emph{Direct Sampling Methods for Inverse Interface Problems}.'' The submitted manuscript
contains a concise proof, while this supplementary note records the intermediate
rational inequalities, radial monotonicity estimates, and polynomial sign
checks. Standard power-series identities for the logarithm and the dilogarithm
are used below; see
\cite[Sec.~25.12]{NIST:DLMF}. Polynomial sign checks are certified by Sturm's
theorem; see \cite[Sec.~2.2.2]{basu2006algorithms}.
\end{abstract}
\section{Full algebraic verification for the almost orthogonal property}
When $\theta_x = \theta_z$, the sign-normalized kernel used in the localization analysis is given by
\begin{equation}
    K(x,z) = d\frac{\sum_{n=1}^{\infty}\frac{n^{2}}{n+\alpha}r_{x}^{n}r_{z}^{n}}{\left(\frac{r_{x}^{2}(1+r_{x}^{2})}{(1-r_{x}^{2})^{3}}\right)^{1/4}\left(\sum_{n\in\mathbb{Z}}\frac{n^{2}}{(|n|+\alpha)^{2}}r_{x}^{2|n|}\right)^{1/4}}.
\end{equation}
where $d>0$. This sign normalization does not affect peak locations, and $K(x,z)^4$ is proportional to the following function.

\begin{theorem}
Let $$f(x) = \frac{\left(\sum_{n=1}^{\infty}\frac{n^{2}}{n+\alpha}c^{n}x^{n}\right)^{4}}{\left(\frac{x^{2}(1+x^{2})}{(1-x^{2})^{3}}\right)\left(\sum_{n=1}^{\infty}\frac{n^{2}}{(n+\alpha)^{2}}x^{2n}\right)}.$$
Furthermore, suppose $c>0.005$ and $\alpha=1$. Then $f'(x)<0$
for all $c\leq x<1$. Hence $f$ is strictly decreasing on $[c,1)$.
\end{theorem}

\begin{proof}
Let 
$$S_1(x) = \sum_{n=1}^{\infty}\frac{n^{2}}{n+1}c^{n}x^{n},\quad S_2(x) = \sum_{n=1}^{\infty}\frac{n^{2}}{(n+1)^{2}}x^{2n}.$$
Then we obtain 
$$f(x) = \frac{(1-x^2)^{3}S_1^{4}(x)}{x^{2}(1+x^{2})S_2(x)}.
$$
It suffices to show that for $c\in(0.005, 1)$,
\begin{equation}
    \left(\log f(x)\right)^{'} = 4\frac{S_1^{'}(x)}{S_1(x)}-\frac{S_2^{'}(x)}{S_2(x)}-\frac{6x}{1-x^2}-\frac{2}{x}-\frac{2x}{1+x^2}<0, \quad c<x<1.
    \label{goal}
\end{equation}
By direct computation,
$$S_1(x) = \sum_{n=1}^{\infty}\frac{n^{2}}{n+1}c^{n}x^{n} = \sum_{n=1}^{\infty}\left(n-1+\frac{1}{n+1}\right)c^{n}x^{n} = \frac{c^2 x^2}{(1-cx)^2}-\frac{\log(1-cx)}{cx}-1,
$$
$$S_1^{'}(x) = \frac{2c^2 x}{(1-cx)^{3}}+\frac{1}{x(1-cx)}+\frac{\log(1-cx)}{cx^2},
$$
$$S_2(x) = \sum_{n=1}^{\infty}x^{2n}-2\sum_{n=1}^{\infty}\frac{x^{2n}}{n+1}+\sum_{n=1}^{\infty}\frac{x^{2n}}{(n+1)^2}=\frac{1}{1-x^2}+\frac{2\log(1-x^2)}{x^2}+\frac{\text{Li}_2(x^2)}{x^2},
$$
$$S_2^{'}(x) = \frac{2x}{(1-x^2)^2}-\frac{4}{x(1-x^2)}-\frac{6\log(1-x^2)}{x^3}-\frac{2\text{Li}_2(x^2)}{x^3}.
$$
To achieve our goal, we first demonstrate (\ref{goal}) holds for $x=c$ and $c\in(0.005, 0.27)$. we prove the following inequalities (\ref{S_1_log_ine})(\ref{S_1_poly_ine})(\ref{S_2_ine}):
\begin{equation}
    \frac{\frac{1}{c(1-c^2)}+\frac{\log(1-c^2)}{c^3}}{\frac{c^4}{(1-c^2)^{2}}-\frac{\log(1-c^2)}{c^2}-1}<-\frac{7c}{8(1-c^2)}-\frac{3c}{8(1+c^2)}+\frac{1}{c}, \quad c\in(0.005, 0.27).\label{S_1_log_ine}
\end{equation}
\begin{equation}
    \frac{\frac{2c^3}{(1-c^2)^3}}{\frac{c^4}{(1-c^2)^{2}}-\frac{\log(1-c^2)}{c^2}-1}<\frac{21c}{8(1-c^2)}+\frac{5c}{4(1+c^2)}+\frac{1}{8}c, \quad c\in(0.005, 0.27).\label{S_1_poly_ine}
\end{equation}
\begin{equation}
    \frac{\frac{2c}{(1-c^2)^2}-\frac{4}{c(1-c^2)}-\frac{6\log(1-c^2)}{c^3}-\frac{2\text{Li}_2(c^2)}{c^3}}{\frac{1}{1-c^2}+\frac{2\log(1-c^2)}{c^2}+\frac{\text{Li}_2(c^2)}{c^2}}>\frac{2}{c}+\frac{c}{1-c^2}+\frac{3c}{2(1+c^2)}+\frac{c}{2}, \quad c\in(0.005, 1).\label{S_2_ine}
\end{equation}
To prove (\ref{S_1_log_ine}), note that (\ref{S_1_log_ine}) is equivalent to the following
\begin{equation}
    (-10c^8+15c^6+8c^4-21c^2+8)\log(1 - c^2)+ (  16c^8 - 17c^4+8c^2)<0
\end{equation}
Let $t = c^2$. It suffices to show that
\begin{equation}
    \log(1-t)+\frac{16t^4 - 17t^2+8t}{-10t^4+15t^3+8t^2-21t+8}<0,\quad t\in(2.5\times10^{-5}, 0.0729).\label{S_1_log_ine_1}
\end{equation}
By direct computation, the  inequality (\ref{S_1_log_ine_1}) holds for $t = 2.5\times10^{-5}$ and $t=0.0729$. Let $L_1(t)$ denote the left-hand side of (\ref{S_1_log_ine_1}). Differentiating $L_{1}(t)$ yields:
$$L_1^{'}(t)=\frac{t^2(100t^4+140t^3-79t^2-67t+4)}{(t - 1)(- 10t^3 + 5t^2 + 13t - 8)^2}
$$
Note that the denominator of $L_1^{'}(t)$ is strictly negative for $t\in(2.5\times10^{-5}, 0.0729)$. For the numerator, the part $t^2>0$, and the part $100t^4+140t^3-79t^2-67t+4$ is positive for $t=2.5\times10^{-5}$, and negative for $t=0.0729$. Additionally, the part $100t^4+140t^3-79t^2-67t+4$ is strictly decreasing on the interval $(2.5\times10^{-5}, 0.0729)$. Therefore, we conclude that $L_1(t)$ first strictly decreases from $2.5\times10^{-5}$ to some point in $(2.5\times10^{-5}, 0.0729)$, and then increases up to 0.0729. Thus, we obtain the desired $L_1(t)<0$, which proves (\ref{S_1_log_ine_1}), and consequently (\ref{S_1_log_ine}).

To prove (\ref{S_1_poly_ine}), we note that it is equivalent to the following inequality for $c\in(0.005, 0.27)$
$$(c^8-13c^6-9c^4+53c^2-32)\log(1-c^2)+(-32c^2+37c^4+7c^6-2c^8)>0.
$$
Let $t = c^2$. It suffices to show that
\begin{equation}
\log(1-t)<\frac{32t-37t^2-7t^3+2t^4}{t^4 -13t^3-9t^2+53t-32},\quad t\in(2.5\times10^{-5}, 0.0729).\label{S_1_poly_subine}    
\end{equation}
Note that $\log(1-t)<-t-t^2/2$. Thus, for $t\in(2.5\times10^{-5}, 0.0729)$, we have:
$$(-t-\frac{t^2}{2})(t^4 -13t^3-9t^2+53t-32)-(32t-37t^2-7t^3+2t^4) = -\frac{1}{2}t^3(t^3+2t^2-18t-5)>0.
$$
Therefore, the desired inequality (\ref{S_1_poly_subine}) holds, which completes the proof of(\ref{S_1_poly_ine}).

To prove (\ref{S_2_ine}), we note that it is equivalent to the following for $c\in(0.005, 1)$:
$$(20-8c^2-34c^4+20c^6 +2c^8)\log(1-c^2) + (8-2c^2 - 15c^4+8c^6+c^8)\text{Li}_2(c^2)<c^8+17c^6-2c^4-12c^2.
$$
Let $t = c^2$. Note that
$$8-2t - 15t^2+8t^3+t^4 = (t-1)(t^3 + 9t^2 - 6t - 8)>0,\quad t\in(2.5\times 10^{-5}, 1),
$$
and
$$20-8t-34t^2+20t^3 +2t^4=2(t-1)(t^3 + 11t^2 - 6t - 10)>0,\quad t\in(2.5\times 10^{-5}, 1).
$$
We have the following:
$$\text{Li}_{2}(t) = \sum_{k=1}^{\infty}\frac{t^k}{k^2}\leq t+\frac{t^2}{4}+\sum_{k=3}^{\infty}\frac{t^k}{3^2}=t+\frac{t^2}{4}+\frac{t^3}{9(1-t)},\quad \log(1-t)\leq -t-\frac{t^2}{2}-\frac{t^3}{3}-\frac{t^4}{4}.
$$
Therefore, to prove (\ref{S_2_ine}), it suffices to show for $t\in(2.5\times 10^{-5}, 1)$:
$$(20-8t-34t^2+20t^3 +2t^4)(-t-\frac{t^2}{2}-\frac{t^3}{3}-\frac{t^4}{4}) + (8-2t - 15t^2+8t^3+t^4)(t+\frac{t^2}{4}+\frac{t^3}{9(1-t)})<t^4+17t^3-2t^2-12t.
$$
This is equivalent to:
\begin{equation}
    \frac{- 18t^8 - 204t^7 + 35t^6 + 120t^5 - 51t^4 - 10t^3}{36}<0.\label{S2-poly-bound}
\end{equation}
Note that:
$$- 18t^8 - 204t^7 + 35t^6 + 120t^5 - 51t^4 - 10t^3 = -t^3(18t^5 + 204t^4 - 35t^3 - 120t^2 + 51t + 10).
$$
Let
\[
U(t)=18t^5+204t^4-35t^3-120t^2+51t+10.
\]
By Sturm's theorem, \(U\) has no zero in \((2.5\times10^{-5},1)\).
Since \(U(2.5\times10^{-5})>0\), it follows that \(U(t)>0\) on this
interval. Hence \eqref{S2-poly-bound} holds. This verifies that (\ref{S_2_ine}) holds for $c\in(0.005, 1)$. At this stage, we have proved for $c\in(0.005, 0.27)$,
$$(\log f(x))^{'}\bigg|_{x=c}<0.
$$
Next, We prove that for $c\in(0.27, 1)$, the following also holds:
$$(\log f(x))^{'}\bigg|_{x=c}<0.$$
We prove the following inequalities (\ref{S_1_ine_2})(\ref{S_2_ine_2}):
\begin{equation}
    \frac{\frac{1}{c(1-c^2)}+\frac{\log(1-c^2)}{c^3}+\frac{2c^3}{(1-c^2)^3}}{\frac{c^4}{(1-c^2)^{2}}-\frac{\log(1-c^2)}{c^2}-1}<\frac{2c}{1-c^2}+\frac{c}{2(1+c^2)}+\frac{1}{c}+\frac{5}{4}c-\frac{c}{1+c^4}, \quad c\in(0.27, 1).\label{S_1_ine_2}
\end{equation}
\begin{equation}
    \frac{\frac{2c}{(1-c^2)^2}-\frac{4}{c(1-c^2)}-\frac{6\log(1-c^2)}{c^3}-\frac{2\text{Li}_2(c^2)}{c^3}}{\frac{1}{1-c^2}+\frac{2\log(1-c^2)}{c^2}+\frac{\text{Li}_2(c^2)}{c^2}}>\frac{2}{c}+\frac{2c}{1-c^2}-\frac{4c}{1+c^4}+5c, \quad c\in( 0.27, 1).\label{S_2_ine_2}
\end{equation}
To prove (\ref{S_1_ine_2}), we note that it is equivalent to the following for $c\in( 0.27, 1)$,
\begin{equation}
    \frac{(-8-11c^2-6c^4-14c^6+2c^8+5c^{10})\log(1-c^2)+(- 8c^2 - 15c^4 - 14c^6 - 23c^8 - 10c^{10})}{4c^3(c^2 - 1)(c^2 + 1)(c^4 + 1)}<0.
\end{equation}
Note that for $c\in( 0.27, 1)$,
$$-8-11c^2-6c^4-14c^6+2c^8+5c^{10}<0.
$$
Therefore, it suffices to show that for $c\in( 0.27, 1)$,
\begin{equation}
        \log(1-c^2)+\frac{- 8c^2 - 15c^4 - 14c^6 - 23c^8 - 10c^{10}}{-8-11c^2-6c^4-14c^6+2c^8+5c^{10}}<0.
    \label{stage2-S1-expanded}
\end{equation}
Let $t=c^2$. and let $T_{1}(t)$ denote the left-hand side of (\ref{stage2-S1-expanded}). Computing the derivative of $T_{1}$ yields:
$$T_1^{'}(t) = \frac{t^2(t+1)(25t^7 + 90t^6 + 99t^5 + 120t^4 + 123t^3 + 66t^2 - 15t + 4)}{(t - 1)(- 5t^5 - 2t^4 + 14t^3 + 6t^2 + 11t + 8)^2
}<0.
$$
Therefore, the left-hand side of (\ref{stage2-S1-expanded}) is decreasing on $(0.27, 1)$. Since
(\ref{stage2-S1-expanded}) holds when $c=0.27$, the inequality holds for $c\in( 0.27, 1)$. Hence, (\ref{S_1_ine_2}) is true.

To prove (\ref{S_2_ine_2}), we note that it is equivalent to the following for $c\in(0.27,1)$:
\begin{equation}
\begin{aligned}
\frac{(-2+10c^2-12c^4+20c^6-26c^8+10c^{10})\log(1-c^2)}{c^3(c^2 - 1)^2(c^4 + 1)}\\
    +\frac{(3c^2-4c^4+8c^6-12c^8+5c^{10})\text{Li}_{2}(c^2)- 2c^2 + 7c^4 - 3c^6 + 11c^8 - 5c^{10}}{c^3(c^2 - 1)^2(c^4 + 1)}>0.
\end{aligned}
\end{equation}
By letting $t=c^2$, it suffices to show that for $t\in( 0.0729, 1)$,
$$(-2+10t-12t^2+20t^3-26t^4+10t^{5})\log(1-t)+(3t-4t^2+8t^3-12t^4+5t^{5})\text{Li}_{2}(t)- 2t + 7t^2 - 3t^3 + 11t^4 - 5t^{5}>0.
$$
Note that we have the following:
$$3t-4t^2+8t^3-12t^4+5t^{5} = t(t-1)(5t^3-7t^2+t-3)>0,\quad t\in( 0.0729, 1),
$$
and
$$\text{Li}_{2}(t) = \sum_{k=1}^{\infty}\frac{t^k}{k^2}\geq t.
$$
We only need to show that for $t\in( 0.0729, 1)$,
\begin{equation}
    (-2+10t-12t^2+20t^3-26t^4+10t^{5})\log(1-t)+(3t-4t^2+8t^3-12t^4+5t^{5})t- 2t + 7t^2 - 3t^3 + 11t^4 - 5t^{5}>0,
\end{equation}
which is equivalent to:
\begin{equation}
    2(t-1)(5t^4 - 8t^3 + 2t^2 - 4t + 1)\log(1-t)+t(5t^5 - 17t^4 + 19t^3 - 7t^2 + 10t - 2)>0,\quad t\in( 0.0729, 1).\label{stage2_last}
\end{equation}
Let $q(t)=5t^4-8t^3+2t^2-4t+1$. By Sturm's theorem,
the number of sign variations of the Sturm sequence of $q$ is $3$ at
$t=0.0729$, $2$ at $t=0.255$, and $2$ at $t=1$. Hence $q$ has exactly
one zero in $(0.0729,0.255)$ and no zero in $[0.255,1)$. Since
$q(0.0729)>0$ and $q(0.255)<0$, the condition $q(t)>0$ implies
$0.0729<t<0.255$. Therefore, when $q(t)>0$, we have
$$2(t-1)(5t^4 - 8t^3 + 2t^2 - 4t + 1)<0,\quad\log(1-t)<-t.
$$
Therefore, we have, for $5t^4 - 8t^3 + 2t^2 - 4t + 1>0$,
\begin{equation}
    \begin{aligned}
        2(t-1)(5t^4 - 8t^3 + 2t^2 - 4t + 1)\log(1-t)&+t(5t^5 - 17t^4 + 19t^3 - 7t^2 + 10t - 2)\\
        >2(t-1)(5t^4 - 8t^3 + 2t^2 - 4t + 1)(-t)+t(5t^5 - 17t^4 + 19t^3 -& 7t^2 + 10t - 2)=-t^3(5t^3 - 9t^2 + t - 5)>0.
    \end{aligned}
\end{equation}
Moreover, when $5t^4 - 8t^3 + 2t^2 - 4t + 1\leq0$, we have the following:
$$2(5t^4 - 8t^3 + 2t^2 - 4t + 1)<0,\quad\text{and}\quad(t-1)\log(1-t)=t-\sum_{k=2}^{\infty}\frac{t^k}{k(k-1)}<t-\frac{t^2}{2}.
$$
Therefore, we have, for $5t^4 - 8t^3 + 2t^2 - 4t + 1\leq0$,
\begin{equation}
    \begin{aligned}
        2(t-1)(5t^4 - 8t^3 + 2t^2 - 4t + 1)\log(1-t)&+t(5t^5 - 17t^4 + 19t^3 - 7t^2 + 10t - 2)\\
        >(5t^4 - 8t^3 + 2t^2 - 4t + 1)(2t-t^2)+t(5t^5 &- 17t^4 + 19t^3 - 7t^2 + 10t - 2)=t^2(t^2+1)(t+1)>0.
    \end{aligned}
\end{equation}
Therefore, (\ref{stage2_last}) is proved, which implies that (\ref{S_2_ine_2}) is true. Combined with (\ref{S_1_ine_2}), we have successfully proved that for $c\in(0.27, 1)$, we also have
$$(\log f(x))^{'}\bigg|_{x=c}<0.$$
At this stage, we have established that
$$(\log f(x))^{'}\bigg|_{x=c}<0
$$
holds for any $c\in(0.005,1)$.

Next, we show that for $c\in(0.005,1)$, 
$$(\log f(x))^{'}<0\quad\text{for all}\quad x>c.
$$
To this end, it suffices to show that:
\begin{equation}
    \frac{S_2^{'}(x)}{S_2(x)}>4\frac{S_1^{'}(x)}{S_1(x)}-\left(\frac{6x}{1-x^2}+\frac{2}{x}+\frac{2x}{1+x^2}\right).
\end{equation}
We show the following two inequalities (\ref{interval1_extend})(\ref{interval2_extend}).

For $0.005<c<0.27$ and $c<x<\min\{\frac{0.27^2}{c},1\}$, by (\ref{S_2_ine}), we only need to show that:
\begin{equation}
    \frac{2}{x}+\frac{x}{1-x^2}+\frac{3x}{2(1+x^2)}+\frac{x}{2}>4\frac{S_1^{'}(x)}{S_1(x)}-\left(\frac{6x}{1-x^2}+\frac{2}{x}+\frac{2x}{1+x^2}\right). \label{interval1_extend}
\end{equation}
To show (\ref{interval1_extend}), since $\sqrt{cx}\in(0.005, 0.27)$, we use (\ref{S_1_log_ine})(\ref{S_1_poly_ine}) to get
\begin{equation}
    4\frac{\frac{1}{x(1-cx)}+\frac{\log(1-cx)}{cx^2}+\frac{2c^2x}{(1-cx)^3}}{\frac{c^2x^2}{(1-cx)^{2}}-\frac{\log(1-cx)}{cx}-1}\sqrt{\frac{x}{c}}<\frac{7\sqrt{cx}}{1-cx}+\frac{7\sqrt{cx}}{2(1+cx)}+\frac{4}{\sqrt{cx}}+\frac{1}{2}\sqrt{cx}.
\end{equation}
Therefore, we have:
\begin{equation}
    4\frac{S_1^{'}(x)}{S_1(x)}<\sqrt{\frac{c}{x}}\left(\frac{7\sqrt{cx}}{1-cx}+\frac{7\sqrt{cx}}{2(1+cx)}+\frac{4}{\sqrt{cx}}+\frac{1}{2}\sqrt{cx}\right) = \frac{7c}{1-cx}+\frac{7c}{2(1+cx)}+\frac{4}{x}+\frac{1}{2}c.
\label{interval1_extend_mid}
\end{equation}
Comparing (\ref{interval1_extend_mid}) and (\ref{interval1_extend}),
it suffices to show that for $0.005<c<0.27$ and $c<x<\min\{\frac{0.27^2}{c},1\}$,
$$\frac{2}{x}+\frac{x}{1-x^2}+\frac{3x}{2(1+x^2)}+\frac{x}{2}+\frac{6x}{1-x^2}+\frac{2}{x}+\frac{2x}{1+x^2}>\frac{7c}{1-cx}+\frac{7c}{2(1+cx)}+\frac{4}{x}+\frac{1}{2}c.
$$
This is equivalent to
$$\frac{7(x-c)}{(1-x^2)(1-cx)}+\frac{7}{2}\frac{x-c}{(1+x^2)(1+cx)}+\frac{1}{2}(x-c)>0,
$$
which is clearly true. Therefore, (\ref{interval1_extend}) is proved for $0.005<c<0.27$ and $c<x<\min\{\frac{0.27^2}{c},1\}$.

For $0.27<c<x<1$, by (\ref{S_2_ine_2}), we only need to show that:
\begin{equation}
    \frac{2}{x}+\frac{2x}{1-x^2}-\frac{4x}{1+x^4}+5x>4\frac{S_1^{'}(x)}{S_1(x)}-\left(\frac{6x}{1-x^2}+\frac{2}{x}+\frac{2x}{1+x^2}\right). \label{interval2_extend}
\end{equation}
To show (\ref{interval2_extend}), since $\sqrt{cx}\in(0.27, 1)$, we use (\ref{S_1_ine_2}) to get:
\begin{equation}
        4\frac{\frac{1}{x(1-cx)}+\frac{\log(1-cx)}{cx^2}+\frac{2c^2x}{(1-cx)^3}}{\frac{c^2x^2}{(1-cx)^{2}}-\frac{\log(1-cx)}{cx}-1}\sqrt{\frac{x}{c}}<\frac{8\sqrt{cx}}{1-cx}+\frac{2\sqrt{cx}}{1+cx}+\frac{4}{\sqrt{cx}}+5\sqrt{cx}-\frac{4\sqrt{cx}}{1+c^2x^2}.
\end{equation}
Therefore, we have:
\begin{equation}
    4\frac{S_1^{'}(x)}{S_1(x)}<\sqrt{\frac{c}{x}}\left(\frac{8\sqrt{cx}}{1-cx}+\frac{2\sqrt{cx}}{1+cx}+\frac{4}{\sqrt{cx}}+5\sqrt{cx}-\frac{4\sqrt{cx}}{1+c^2x^2}\right) = \frac{8c}{1-cx}+\frac{2c}{1+cx}+\frac{4}{x}+5c-\frac{4c}{1+c^2x^2}.
\label{interval2_extend_mid}
\end{equation}
Comparing (\ref{interval2_extend_mid}) and (\ref{interval2_extend}),
it suffices to show that for $0.27<c<x<1$,
$$\frac{2}{x}+\frac{2x}{1-x^2}-\frac{4x}{1+x^4}+5x+\frac{6x}{1-x^2}+\frac{2}{x}+\frac{2x}{1+x^2}>\frac{8c}{1-cx}+\frac{2c}{1+cx}+\frac{4}{x}+5c-\frac{4c}{1+c^2x^2}
$$
This is equivalent to
$$\frac{8(x-c)}{(1-cx)(1-x^2)}+\frac{4(cx^3-1)(x-c)}{(1+x^4)(1+c^2x^2)}+5(x-c)+\frac{2(x-c)}{(1+cx)(1+x^2)}>0,
$$
which is true since $\frac{4(cx^3-1)}{(1+x^4)(1+c^2x^2)}>-4$. Therefore, (\ref{interval2_extend}) is proved for $0.27<c<x<1$.

For $0.005<c<0.27$ and $\frac{0.27^2}{c}\leq x<1$, we can obtain the desired results using the same techniques as in the previous case where $0.27<c<x<1$.

Through our previous calculation process, we can get either
$$(\log f(x))^{'} <(c-x)\left(\frac{7}{(1-x^2)(1-cx)}+\frac{7}{2}\frac{1}{(1+x^2)(1+cx)}+\frac{1}{2}\right)<0,
$$
or
$$(\log f(x))^{'} <(c-x)\left(\frac{8}{(1-cx)(1-x^2)}+\frac{4(cx^3-1)}{(1+x^4)(1+c^2x^2)}+5+\frac{2}{(1+cx)(1+x^2)}\right)<0.
$$
Therefore, for $0.005<c<1$, $f$ is strictly decreasing on $[c,1)$.
This completes the proof of Theorem 1.

\end{proof}

\begin{theorem}
Let $$f(x) = \frac{\left(\sum_{n=1}^{\infty}\frac{n^{2}}{n+\alpha}c^{n}x^{n}\right)^{4}}{\left(\frac{x^{2}(1+x^{2})}{(1-x^{2})^{3}}\right)\left(\sum_{n=1}^{\infty}\frac{n^{2}}{(n+\alpha)^{2}}x^{2n}\right)}.$$
Suppose $\alpha=1$ and $0.005<c<1$. Then the following hold.

If $0.005<c<0.5$, then
\[
(\log f(x))'>0
\quad\text{for all}\quad
0.0045<x<0.9c.
\]

If $0.5\leq c<1$, then
\[
(\log f(x))'>0
\quad\text{for all}\quad
0.0045<x<0.95c.
\]
Equivalently, since $f(x)>0$, the function $f$ is strictly increasing on
the corresponding intervals.
\end{theorem}
\begin{proof}
We use the same notation as in the proof of Theorem 1. It suffices to show that for $0.005<c<0.5$:
\begin{equation}
    (\log f(x))^{'} = 4\frac{S_1^{'}(x)}{S_1(x)}-\frac{S_2^{'}(x)}{S_2(x)}-\frac{6x}{1-x^2}-\frac{2}{x}-\frac{2x}{1+x^2}>0,\quad 0.0045<x\leq 0.9c,\label{goal1}
\end{equation}
and for $0.5\leq c<1$:
\begin{equation}
    (\log f(x))^{'} = 4\frac{S_1^{'}(x)}{S_1(x)}-\frac{S_2^{'}(x)}{S_2(x)}-\frac{6x}{1-x^2}-\frac{2}{x}-\frac{2x}{1+x^2}>0,\quad 0.0045<x\leq 0.95c.\label{goal2}
\end{equation}
To achieve this, we first demonstrate that (\ref{goal1}) holds for $x=0.9c$ and $0.005<c<0.5$, and (\ref{goal2}) holds for $x=0.95c$ and $0.5\leq c<1$.

To prove that $(\log f(x))^{\prime}|_{x=0.95c}>0$ for $0.5\le c<1$, we prove the following inequalities (\ref{goal1_S1_ine})(\ref{goal1_S2_ine}).

Let $\mu = \frac{20}{19}$, then $c=\mu x$. We show that:
\begin{equation}
    \frac{\frac{1}{x(1-\mu x^2)}+\frac{\log(1-\mu x^2)}{\mu x^3}+\frac{2\mu^2x^3}{(1-\mu x^2)^3}}{\frac{\mu^2x^4}{(1-\mu x^2)^{2}}-\frac{\log(1-\mu x^2)}{\mu x^2}-1}>\frac{11x}{4(1-x^2)}+\frac{2x}{1+x^2}+\frac{1}{x}-\frac{7}{4}x, \quad x\in[0.48/\mu, 1/\mu).\label{goal1_S1_ine}
\end{equation}
\begin{equation}
    \frac{\frac{2x}{(1-x^2)^2}-\frac{4}{x(1-x^2)}-\frac{6\log(1-x^2)}{x^3}-\frac{2\text{Li}_2(x^2)}{x^3}}{\frac{1}{1-x^2}+\frac{2\log(1-x^2)}{x^2}+\frac{\text{Li}_2(x^2)}{x^2}}<\frac{2}{x}+\frac{5x}{1-x^2}+\frac{6x}{1+x^2}-7x, \quad x\in(0.005/\mu, 1/\mu).\label{goal1_S2_ine}
\end{equation}
To prove (\ref{goal1_S1_ine}), let $y = \mu x^2$. This is equivalent to proving the following inequality:
\begin{equation}
\frac{\frac{1}{1-y}+\frac{\log(1-y)}{y}+\frac{2y^2}{(1-y)^3}}{\frac{y^2}{(1-y)^{2}}-\frac{\log(1-y)}{y}-1}>\frac{11y/\mu}{4(1-y/\mu)}+\frac{2y/\mu}{1+y/\mu}+1-\frac{7}{4}y/\mu, \quad y\in[0.48^2/\mu, 1/\mu).\label{S_1y_ine}
\end{equation}
We find the common denominator on both sides of the inequality, and it suffices to prove:
\begin{equation}
\begin{aligned}
    \frac{(- 48013y^6 + 180139y^5 - 343539y^4 + 365913y^3 - 117700y^2 - 100800y + 64000)\log(1-y)}{80y(19y - 20)(19y + 20)(y - 1)^3}\\
    +\frac{64000y- 68800y^2 - 149700y^3 + 309833y^4 - 245119y^5 + 96026y^6}{80y(19y - 20)(19y + 20)(y - 1)^3}>0
\end{aligned}
\end{equation}
Note that:
\begin{align*}
- 48013y^6 + 180139y^5 - 343539y^4 + 365913y^3 - 117700y^2 - 100800y + 64000\\
= -(y-1)^3(48013y^3 - 36100y^2 + 91200y + 64000)>0.
\end{align*}
Denoting $R(y) = - 48013y^6 + 180139y^5 - 343539y^4 + 365913y^3 - 117700y^2 - 100800y + 64000$, and $P(y) = 64000y- 68800y^2 - 149700y^3 + 309833y^4 - 245119y^5 + 96026y^6$,
we need to show that for $y\in[0.48^2/\mu, 1/\mu)$
$$\log(1-y)>-\frac{P(y)}{R(y)}.
$$
We first check that it is true for $y = 0.48^2/\mu=0.21888$. Then, taking derivatives for both sides, it suffices to show that
$$-\frac{1}{1-y}+\frac{P^{'}R-PR^{'}}{R^2}>0.
$$
A direct simplification gives
\[
-\frac{1}{1-y}+\frac{P'R-PR'}{R^2}
=
\frac{y^2 p(y)}{(1-y)^4(48013y^3-36100y^2+91200y+64000)^2},
\]
where
\[
\begin{aligned}
p(y)=&2305248169y^7-4853154040y^6+2209071271y^5
+13054323160y^4\\
&-7282525200y^3-6421088000y^2+4679680000y-563200000.
\end{aligned}
\]
By Sturm's theorem, the number of sign variations of the Sturm sequence
of $p$ is $3$ at $y=684/3125$ and also $3$ at $y=19/20$. Hence $p$
has no zero on $[684/3125,19/20]$. Since $p(684/3125)>0$, we have
$p(y)>0$ throughout this interval. Therefore, we have
$$-\frac{1}{1-y}+\frac{P^{'}R-PR^{'}}{R^2}>0.$$
Thus, (\ref{goal1_S1_ine}) is proved.

To prove (\ref{goal1_S2_ine}), it is enough to prove the following
inequality for $0<x<1$:
\begin{equation}
    \frac{(10-2x^2-20x^4+26x^6-14x^8)\log(1-x^2)+(4-9x^4+12x^6-7x^8)\text{Li}_2(x^2)+(6x^2 + 2x^4 - 9x^6 + 7x^8)}{x^3(x^2 - 1)^2(x^2 + 1)}>0.\label{whole_st2}
\end{equation}
Let $t=x^2$. Note that
$$10-2t-20t^2+26t^3-14t^4 = -2(t-1)(7t^3 - 6t^2 + 4t + 5)>0,
$$
$$4-9t^2+12t^3-7t^4 = -(t-1)(7t^3 - 5t^2 + 4t + 4)>0,
$$
$$(t-1)\log(1-t)=t-\sum_{k=2}^{\infty}\frac{t^k}{k(k-1)}<t-\frac{t^2}{2}-\frac{t^3}{6}-\frac{t^4}{12}-\frac{t^5}{20}-\frac{t^6}{30}-\frac{t^7}{42},
$$
$$\text{Li}_2(t)>t+\frac{t^2}{4}+\frac{t^3}{9}+\frac{t^4}{16}+\frac{t^5}{25}+\frac{t^6}{36}.
$$
Therefore, to prove (\ref{whole_st2}), it suffices to show:
\begin{align*}
    -2(7t^3 - 6t^2 + 4t + 5)(t-\frac{t^2}{2}-\frac{t^3}{6}-\frac{t^4}{12}-\frac{t^5}{20}-\frac{t^6}{30}-\frac{t^7}{42})\\
    +(4-9t^2+12t^3-7t^4)(t+\frac{t^2}{4}+\frac{t^3}{9}+\frac{t^4}{16}+\frac{t^5}{25}+\frac{t^6}{36})+6x^2 + 2x^4 - 9x^6 + 7x^8>0,
\end{align*}
which is equivalent to for $t\in(0.005^2/\mu^2, 1/\mu^2)$,
$$t^3(3500t^7 + 5904t^6 + 7131t^5 + 17228t^4 + 30205t^3 + 33432t^2 - 21000t + 2800
)>0.
$$
Let
\[
p(t)=3500t^7+5904t^6+7131t^5+17228t^4+30205t^3
+33432t^2-21000t+2800.
\]
By Sturm's theorem, the number of sign variations of the Sturm sequence
of $p$ is $3$ at $t=0$ and $3$ at $t=1$. Hence $p$ has no zero in
$(0,1)$. Since $p(0)=2800>0$, it follows that $p(t)>0$ for $0<t<1$. To this stage, we have proved (\ref{whole_st2}) is true, and (\ref{goal1_S2_ine}) follows. Therefore, for $0.5\leq c<1$,
$$(\log f(x))^{'}\bigg|_{x=0.95c}>0.
$$
Then we show that for $0.005<c<0.5$, 
$$(\log f(x))^{'}\bigg|_{x=0.9c}>0.
$$
Let $\beta = \frac{10}{9}$, then $c=\beta x$. Note that
\[
(0.005/\beta, 0.5/\beta)\subset(0.005/\beta, 0.48/\mu),
\]
since $0.5/\beta=0.45<0.48/\mu=0.456$. We show that:
\begin{equation}
    \frac{\frac{1}{x(1-\beta x^2)}+\frac{\log(1-\beta x^2)}{\beta x^3}+\frac{2\beta^2x^3}{(1-\beta x^2)^3}}{\frac{\beta^2x^4}{(1-\beta x^2)^{2}}-\frac{\log(1-\beta x^2)}{\beta x^2}-1}>\frac{11x}{4(1-x^2)}+\frac{5x}{4(1+x^2)}+\frac{1}{x}-\frac{17}{16}x, \quad x\in(0.005/\beta, 0.48/\mu).\label{goal2_S1_ine}
\end{equation}
\begin{equation}
    \frac{\frac{2x}{(1-x^2)^2}-\frac{4}{x(1-x^2)}-\frac{6\log(1-x^2)}{x^3}-\frac{2\text{Li}_2(x^2)}{x^3}}{\frac{1}{1-x^2}+\frac{2\log(1-x^2)}{x^2}+\frac{\text{Li}_2(x^2)}{x^2}}<\frac{2}{x}+\frac{5x}{1-x^2}+\frac{3x}{1+x^2}-\frac{17}{4}x, \quad x\in(0.005/\beta, 0.48/\mu).\label{goal2_S2_ine}
\end{equation}
To prove (\ref{goal2_S1_ine}), let $w = \beta x^2 = cx$. Then (\ref{goal2_S1_ine}) is equivalent to
\begin{equation}
    \frac{\frac{1}{1-w}+\frac{\log(1-w)}{w}+\frac{2w^2}{(1-w)^3}}{\frac{w^2}{(1-w)^{2}}-\frac{\log(1-w)}{w}-1}>\frac{11w/\beta}{4(1-w/\beta)}+\frac{5w/\beta}{4(1+w/\beta)}+1-\frac{17}{16}w/\beta, \quad w\in[0.005^2/\beta, 0.48^2\beta/\mu^2).
\end{equation}
We find the common denominator on both sides of the inequality, and it suffices to prove
\begin{equation}
\begin{aligned}
    \frac{(- 12393w^6 + 43659w^5 - 98919w^4 + 126733w^3 - 37380w^2 - 53700w + 32000)\log(1-w)}{160w(9w - 10)(9w + 10)(w - 1)^3}\\
    +\frac{32000w- 37700w^2 - 53380w^3 + 103473w^4 - 63099w^5 + 24786w^6}{160w(9w - 10)(9w + 10)(w - 1)^3}>0.
\end{aligned}
\end{equation}
Note that
\begin{align*}
    - 12393w^6 + 43659w^5 - 98919w^4 + 126733w^3 - 37380w^2 - 53700w + 32000\\ = -(w-1)^3(12393w^3 - 6480w^2 + 42300w + 32000)>0.
\end{align*}
$$(w-1)\log(1-w)<w-\frac{w^2}{2}-\frac{w^3}{6}.
$$
Therefore, it suffices to show that
\begin{align*}
    -(w-1)^2(12393w^3 - 6480w^2 + 42300w + 32000)(w-\frac{w^2}{2}-\frac{w^3}{6})\\
    +32000w- 37700w^2 - 53380w^3 + 103473w^4 - 63099w^5 + 24786w^6>0.
\end{align*}
This is equivalent to for $w\in[0.005^2/\beta, 0.48^2\beta/\mu^2) = (2.25\times 10^{-5}, 0.23104)$,
$$\frac{1}{6}w^3(12393w^5 + 5913w^4 + 48213w^3 - 47119w^2 + 15980w + 1100)>0,
$$
which holds clearly. Therefore, (\ref{goal2_S1_ine}) is true. Next, we prove (\ref{goal2_S2_ine}). We note that (\ref{goal2_S2_ine}) is equivalent to for $x\in(0.005/\beta, 0.48/\mu)$,
\begin{equation}
\begin{aligned}
\frac{(40-10x^2-54x^4+58x^6-34x^8)\log(1-x^2)+(16-x^2-23x^4+25x^6-17x^8)\text{Li}_2(x^2)}{4x^3(x^2 - 1)^2(x^2 + 1)}\\
\frac{+24x^2 + 7x^4 - 24x^6 + 17x^8}{4x^3(x^2 - 1)^2(x^2 + 1)}>0.    
\end{aligned}
 \label{goal2_S2_mid}
\end{equation}
Let $t=x^2$. Note that
$$40-10t-54t^2+58t^3-34t^4 = -2(t-1)(17t^3 - 12t^2+ 15t + 20)>0,
$$
$$16-t-23t^2+25t^3-17t^4 = -(t-1)(17t^3 - 8t^2 + 15t + 16)>0,
$$
$$(t-1)\log(1-t)<t-\frac{t^2}{2}-\frac{t^3}{6}-\frac{t^4}{12},\quad \text{Li}_2(t)>t+\frac{t^2}{4}+\frac{t^3}{9}+\frac{t^4}{16}.
$$
Therefore, it suffices to show that
\begin{align*}
    -2(17t^3 - 12t^2+ 15t + 20)(t-\frac{t^2}{2}-\frac{t^3}{6}-\frac{t^4}{12})+(16-t-23t^2+25t^3-17t^4)(t+\frac{t^2}{4}+\frac{t^3}{9}+\frac{t^4}{16})\\
    +(24t + 7t^2 - 24t^3 + 17t^4)>0.
\end{align*}
This is equivalent to show for $t\in(0.005^2/\beta^2, 0.48^2/\mu^2)=(2.025\times 10^{-5}, 0.207936)$,
$$-\frac{1}{144}t^3(153t^5 - 361t^4 - 109t^3 - 307t^2 + 76t - 28)>0,
$$
which holds clearly. Therefore, (\ref{goal2_S2_mid}) is true and we have proved (\ref{goal2_S2_ine}). To this stage, we have shown that for $c\in(0.005, 0.5)$,
$$(\log f(x))^{'}\bigg|_{x=0.9c} > 0.
$$
We now extend the endpoint inequalities to the full intervals. For
$0.5\leq c<1$, we split
\[
0.0045<x<0.95c
\]
into
\[
0.0045<x\leq \frac{0.48^2}{c\mu}
\quad\text{and}\quad
\frac{0.48^2}{c\mu}<x<\frac{c}{\mu}=0.95c.
\]
For $0.005<c<0.5$, we prove the result on
\[
0.0045<x<\frac{c}{\beta}=0.9c.
\]
Specifically, we show that
\begin{equation}
    \frac{S_2^{'}(x)}{S_2(x)}<4\frac{S_1^{'}(x)}{S_1(x)}-\left(\frac{6x}{1-x^2}+\frac{2}{x}+\frac{2x}{1+x^2}\right).
\end{equation}
We show the following two inequalities (\ref{goal2_interval1_extend})(\ref{goal2_interval2_extend}).

Let $\mu = 20/19$. For $0.5\leq c<1$ and $\frac{0.48^2}{c\mu}<x<c/\mu$, by (\ref{goal1_S2_ine}), we only need to show that
\begin{equation}
    \frac{2}{x}+\frac{5x}{1-x^2}+\frac{6x}{1+x^2}-7x<4\frac{S_1^{'}(x)}{S_1(x)}-\left(\frac{6x}{1-x^2}+\frac{2}{x}+\frac{2x}{1+x^2}\right). \label{goal2_interval1_extend}
\end{equation}
To show (\ref{goal2_interval1_extend}), since $cx\in(0.48^2/\mu, 1/\mu)$, we use (\ref{S_1y_ine}) to get
\begin{equation}
    \frac{\frac{1}{1-cx}+\frac{\log(1-cx)}{cx}+\frac{2c^2x^2}{(1-cx)^3}}{\frac{c^2x^2}{(1-cx)^{2}}-\frac{\log(1-cx)}{cx}-1}>\frac{11cx/\mu}{4(1-cx/\mu)}+\frac{2cx/\mu}{1+cx/\mu}+1-\frac{7}{4}cx/\mu.
\end{equation}
Therefore, we have
\begin{equation}
    4\frac{S_1^{'}(x)}{S_1(x)}>\frac{4}{x}\left(\frac{11cx/\mu}{4(1-cx/\mu)}+\frac{2cx/\mu}{1+cx/\mu}+1-\frac{7}{4}cx/\mu\right) = \frac{11c/\mu}{1-cx/\mu}+\frac{8c/\mu}{1+cx/\mu}+\frac{4}{x}-7c/\mu.
\label{goal2_interval1_extend_mid}
\end{equation}
Comparing (\ref{goal2_interval1_extend_mid}) and (\ref{goal2_interval1_extend}),
We only need to show that for $0.5\leq c<1$ and $\frac{0.48^2}{c\mu}<x<c/\mu$,
$$\frac{2}{x}+\frac{5x}{1-x^2}+\frac{6x}{1+x^2}-7x+\frac{6x}{1-x^2}+\frac{2}{x}+\frac{2x}{1+x^2}<\frac{11c/\mu}{1-cx/\mu}+\frac{8c/\mu}{1+cx/\mu}+\frac{4}{x}-7c/\mu.
$$
This is equivalent to
$$\frac{11(x-c/\mu)}{(1-x^2)(1-cx/\mu)}+8\frac{x-c/\mu}{(1+x^2)(1+cx/\mu)}-7(x-c/\mu)<0,
$$
which is clearly true. Therefore, (\ref{goal2_interval1_extend}) is proved for $0.5\leq c<1$ and $\frac{0.48^2}{c\mu}<x<c/\mu$.\\

For $0.005<c<0.5$ and $0.0045<x<0.9c$, by (\ref{goal2_S2_ine}), we only need to show that
\begin{equation}
    \frac{2}{x}+\frac{5x}{1-x^2}+\frac{3x}{1+x^2}-\frac{17}{4}x<4\frac{S_1^{'}(x)}{S_1(x)}-\left(\frac{6x}{1-x^2}+\frac{2}{x}+\frac{2x}{1+x^2}\right). \label{goal2_interval2_extend}
\end{equation}
Let $\beta = 10/9$. To show (\ref{goal2_interval2_extend}), since
\[
cx\in(0.005^2/\beta,0.5^2/\beta)
\subset (0.005^2/\beta,0.48^2\beta/\mu^2),
\]
we use (\ref{goal2_S1_ine}) to obtain
\begin{equation}
        \frac{\frac{1}{1-cx}+\frac{\log(1-cx)}{cx}+\frac{2c^2x^2}{(1-cx)^3}}{\frac{c^2x^2}{(1-cx)^{2}}-\frac{\log(1-cx)}{cx}-1}>\frac{11cx/\beta}{4(1-cx/\beta)}+\frac{5cx/\beta}{4(1+cx/\beta)}+1-\frac{17}{16}cx/\beta.
\end{equation}
Therefore, we have:
\begin{equation}
    4\frac{S_1^{'}(x)}{S_1(x)}>\frac{4}{x}\left(\frac{11cx/\beta}{4(1-cx/\beta)}+\frac{5cx/\beta}{4(1+cx/\beta)}+1-\frac{17}{16}cx/\beta\right) = \frac{11c/\beta}{1-cx/\beta}+\frac{5c/\beta}{1+cx/\beta}+\frac{4}{x}-\frac{17c/\beta}{4}.
\label{goal2_interval2_extend_mid}
\end{equation}
Comparing (\ref{goal2_interval2_extend_mid}) and (\ref{goal2_interval2_extend}),
We only need to show that for $0.005<c<0.5$ and $0.005/\beta<x<c/\beta$,
$$\frac{2}{x}+\frac{5x}{1-x^2}+\frac{3x}{1+x^2}-\frac{17}{4}x+\frac{6x}{1-x^2}+\frac{2}{x}+\frac{2x}{1+x^2}<\frac{11c/\beta}{1-cx/\beta}+\frac{5c/\beta}{1+cx/\beta}+\frac{4}{x}-\frac{17c/\beta}{4}.
$$
This is equivalent to
$$\frac{11(x-c/\beta)}{(1-cx/\beta)(1-x^2)}-\frac{17}{4}(x-c/\beta)+\frac{5(x-c/\beta)}{(1+cx/\beta)(1+x^2)}<0,
$$
which is clearly true. Therefore, (\ref{goal2_interval2_extend}) is proved for $0.005<c<0.5$ and $0.005/\beta<x<c/\beta$.

It remains to handle the lower subinterval for $0.5\leq c<1$. If
\[
0.0045<x\leq \frac{0.48^2}{c\mu},
\]
then
\[
x\leq \frac{0.48^2}{c\mu}\leq \frac{0.48^2}{0.5\mu}=0.43776
<0.45\leq \frac{c}{\beta}.
\]
Moreover,
\[
cx\leq \frac{0.48^2}{\mu}<\frac{0.48^2\beta}{\mu^2}.
\]
Therefore, the same argument used in the proof of
(\ref{goal2_interval2_extend}) applies and gives
\[
(\log f(x))^{'}
>
(c/\beta-x)\left(\frac{11}{(1-x^2)(1-cx/\beta)}
+\frac{5}{(1+x^2)(1+cx/\beta)}-\frac{17}{4}\right)>0.
\]

Combining the estimates above, for $0.005<c<0.5$ and
$0.0045<x<0.9c$, and also for $0.5\leq c<1$ and
$0.0045<x\leq \frac{0.48^2}{c\mu}$, we have
\[
(\log f(x))^{'}
>
(c/\beta-x)\left(\frac{11}{(1-x^2)(1-cx/\beta)}
+\frac{5}{(1+x^2)(1+cx/\beta)}-\frac{17}{4}\right)>0.
\]
For $0.5\leq c<1$ and
$\frac{0.48^2}{c\mu}<x<0.95c=c/\mu$, we have
\[
(\log f(x))^{'}
>
(c/\mu-x)\left(\frac{11}{(1-cx/\mu)(1-x^2)}
+\frac{8}{(1+cx/\mu)(1+x^2)}-7\right)>0.
\]
This completes the proof of Theorem 2.
\end{proof}
\bibliographystyle{siamplain}
\bibliography{References}
\end{document}
